%!TEX TS-program = xelatex
\documentclass[11pt]{article}
\usepackage[margin=1in]{geometry}
\usepackage{fontspec}

\newfontfamily\kurdishfont[Script=Arabic]{Amiri}
% For code listings: FreeMono keeps brackets clear and supports Kurdish text.
\newfontfamily\codemonofont{FreeMono}
\newcommand{\codeleftparen}{{\codemonofont(}}
\newcommand{\coderightparen}{{\codemonofont)}}
\newcommand{\codeleftbracket}{{\codemonofont[}}
\newcommand{\coderightbracket}{{\codemonofont]}}
\newcommand{\codeleftbrace}{{\codemonofont\char`\{}}
\newcommand{\coderightbrace}{{\codemonofont\char`\}}}

\usepackage{hyperref}
\usepackage{booktabs}

% Core packages
\usepackage{amsmath,amssymb}
\usepackage{tikz-cd}
\usepackage{multicol}
\usepackage{xcolor}
\usepackage{listings}

% Code listings
\lstdefinestyle{pythoncode}{
  language=Python,
  basicstyle=\footnotesize\codemonofont,
  keywordstyle=\color{blue!70!black},
  commentstyle=\color{green!40!black},
  stringstyle=\color{orange!50!black},
  showstringspaces=false,
  breaklines=true,
  breakatwhitespace=false,
  columns=fullflexible,
  keepspaces=true,
  tabsize=4,
  frame=single,
  numbers=left,
  numberstyle=\scriptsize\codemonofont\color{gray},
  stepnumber=1,
  numbersep=10pt,
  extendedchars=false,
  literate=
    {(}{{\codeleftparen}}1
    {)}{{\coderightparen}}1
    {[}{{\codeleftbracket}}1
    {]}{{\coderightbracket}}1
    {\{}{{\codeleftbrace}}1
    {\}}{{\coderightbrace}}1
}

\lstdefinestyle{bashcode}{
  language=bash,
  basicstyle=\footnotesize\codemonofont,
  keywordstyle=\color{blue!70!black},
  commentstyle=\color{green!40!black},
  stringstyle=\color{orange!50!black},
  showstringspaces=false,
  breaklines=true,
  breakatwhitespace=false,
  columns=fullflexible,
  keepspaces=true,
  tabsize=4,
  frame=single,
  numbers=left,
  numberstyle=\scriptsize\codemonofont\color{gray},
  stepnumber=1,
  numbersep=10pt,
  extendedchars=false,
  literate=
    {(}{{\codeleftparen}}1
    {)}{{\coderightparen}}1
    {[}{{\codeleftbracket}}1
    {]}{{\coderightbracket}}1
    {\{}{{\codeleftbrace}}1
    {\}}{{\coderightbrace}}1
}

% Paragraphs
\setlength{\parindent}{0pt}
\setlength{\parskip}{1\baselineskip}

\title{Dual-GPU Qwen Knowledge Graph Pipeline}
\author{Data Engineering Team}
\date{\today}

\begin{document}
\maketitle

\section*{Overview}

This document contains a comprehensive toolkit of scripts utilized for generating a medical Knowledge Graph from the Kurdish medical dataset. It leverages a dual-GPU deployment of the Qwen3.6:35b model via Ollama for semantic relation extraction and Neo4j for storage. Below is a summary table containing links to the code for each step in the pipeline, alongside an explanation of its purpose.

\begin{table}[htbp]
\centering
\small
\renewcommand{\arraystretch}{1.3}
\begin{tabular}{@{} l l p{9cm} @{}}
\toprule
\textbf{Script Name} & \textbf{Section Link} & \textbf{Purpose / Why We Use It} \\
\midrule

\multicolumn{3}{@{}l}{\textbf{Data Preparation}} \\
\midrule
Split Dataset & \hyperref[sec:split]{Section \ref*{sec:split}} & Splits the dataset into two parts for parallel processing on two GPUs. \\
\addlinespace

\multicolumn{3}{@{}l}{\textbf{Knowledge Graph Generators}} \\
\midrule
KG Generator (GPU 0) & \hyperref[sec:gengpu0]{Section \ref*{sec:gengpu0}} & Extracts medical triples using Ollama on GPU 0 and stores them in Neo4j. \\
KG Generator (GPU 1) & \hyperref[sec:gengpu1]{Section \ref*{sec:gengpu1}} & Extracts medical triples using Ollama on GPU 1 and stores them in Neo4j. \\
Run KG Script (GPU 0) & \hyperref[sec:runkg0]{Section \ref*{sec:runkg0}} & Shell script to run the GPU 0 knowledge graph generator in the background. \\
Run KG Script (GPU 1) & \hyperref[sec:runkg1]{Section \ref*{sec:runkg1}} & Shell script to run the GPU 1 knowledge graph generator in the background. \\
Stop KG Script & \hyperref[sec:stopkg]{Section \ref*{sec:stopkg}} & Gracefully stops the background Python processes and Ollama servers. \\
\addlinespace

\multicolumn{3}{@{}l}{\textbf{Setup \& Environment}} \\
\midrule
Setup Script & \hyperref[sec:setup]{Section \ref*{sec:setup}} & Prepares the environment by installing dependencies and configuring Neo4j/Ollama. \\
Setup Runner & \hyperref[sec:runsetup]{Section \ref*{sec:runsetup}} & Python wrapper to execute the setup shell script in the background. \\
Model Loader (GPU 0) & \hyperref[sec:loader0]{Section \ref*{sec:loader0}} & A dedicated script to manage loading the Ollama model into GPU 0 memory. \\
Model Loader (GPU 1) & \hyperref[sec:loader1]{Section \ref*{sec:loader1}} & A dedicated script to manage loading the Ollama model into GPU 1 memory. \\
Requirements & \hyperref[sec:requirements]{Section \ref*{sec:requirements}} & Lists the necessary Python packages for the pipeline. \\
\bottomrule
\end{tabular}
\caption{Dual-GPU Knowledge Graph generation scripts categorized by functional domain.}
\end{table}

\clearpage

\section{Split Dataset} \label{sec:split}
\textbf{Purpose:} Splits the dataset into two equal halves to allow parallel processing across the two GPUs.

\begin{lstlisting}[style=pythoncode, caption={split\_dataset.py}]
import json
import math
import os

def split_dataset(input_path, output_part1, output_part2):
    print(f"Reading {input_path}...")
    with open(input_path, 'r', encoding='utf-8') as f:
        data = json.load(f)
    
    total_entries = len(data)
    halfway = math.ceil(total_entries / 2)
    
    part1 = data[:halfway]
    part2 = data[halfway:]
    
    print(f"Total entries: {total_entries}")
    print(f"Part 1: {len(part1)} entries")
    print(f"Part 2: {len(part2)} entries")
    
    with open(output_part1, 'w', encoding='utf-8') as f:
        json.dump(part1, f, ensure_ascii=False, indent=2)
    
    with open(output_part2, 'w', encoding='utf-8') as f:
        json.dump(part2, f, ensure_ascii=False, indent=2)
    
    print(f"Successfully saved to {output_part1} and {output_part2}")

if __name__ == "__main__":
    DATASET_PATH = "kurdish_medical_corpus_kmc.json"
    PART1_PATH = "kurdish_medical_corpus_kmc_part1.json"
    PART2_PATH = "kurdish_medical_corpus_kmc_part2.json"
    
    if os.path.exists(DATASET_PATH):
        split_dataset(DATASET_PATH, PART1_PATH, PART2_PATH)
    else:
        print(f"Error: {DATASET_PATH} not found.")
\end{lstlisting}

\section{Knowledge Graph Generator (GPU 0)} \label{sec:gengpu0}
\textbf{Purpose:} We use this script to extract medical triples from the first half of the Kurdish dataset using the Ollama service on GPU 0.

\begin{lstlisting}[style=pythoncode, caption={kg\_generator\_gpu\_0.py}]
import json
import os
import re
import time
import logging
import threading
from concurrent.futures import ThreadPoolExecutor
from typing import List, Dict, Optional
from tqdm import tqdm
from ollama import Client
from neo4j import GraphDatabase
from dotenv import load_dotenv

# Load environment variables
load_dotenv()

# ─────────────────────────────────────────────────────────
# GPU 0  →  Ollama on port 11434
# ─────────────────────────────────────────────────────────

# Setup Logging
logging.basicConfig(
    level=logging.INFO,
    format='%(asctime)s - %(levelname)s - %(message)s',
    handlers=[
        logging.FileHandler("kg_pipeline_gpu_0.log"),
        logging.StreamHandler()
    ]
)
logger = logging.getLogger(__name__)

# ── Configuration ─────────────────────────────────────────
OLLAMA_HOST   = os.getenv("OLLAMA_HOST_GPU0", "http://localhost:11434")
OLLAMA_MODEL  = "qwen3.6:35b"
NEO4J_URI     = os.getenv("NEO4J_URI",      "bolt://localhost:7687")
NEO4J_USER    = os.getenv("NEO4J_USER",     "neo4j")
NEO4J_PASSWORD= os.getenv("NEO4J_PASSWORD", "password")
DATASET_PATH  = "kurdish_medical_corpus_kmc_part1.json"
CHECKPOINT_PATH = "processed_ids_gpu_0.txt"
OUTPUT_FILE   = "extracted_triples_gpu_0.jsonl"

# Qwen3 generation options – generous token budgets so CoT has room to breathe
OLLAMA_OPTIONS = {
    "temperature": 0.6,      # Qwen3 docs recommend 0.6 for balanced CoT
    "top_p": 0.95,           # Qwen3 recommended value
    "top_k": 20,             # Qwen3 recommended value
    "min_p": 0,
    "num_predict": 8192,     # Large budget: CoT + JSON output
    "num_ctx": 16384,        # Context window – enough for long Kurdish texts + prompt
    "repeat_penalty": 1.0,
}

# ── Robust JSON extractor ──────────────────────────────────
def _extract_json_from_text(text: str) -> Optional[list]:
    """
    Extract the first JSON array from arbitrary text (including CoT reasoning).
    Qwen3 thinks in <think>…</think> blocks before giving the answer;
    we must scan the *entire* response for any valid JSON array.
    """
    if not text:
        return None

    # Strip think blocks if present (Qwen3 CoT)
    text_no_think = re.sub(r'<think>.*?</think>', '', text, flags=re.DOTALL).strip()

    # Prefer text after think block; fall back to full text
    sources = [text_no_think, text] if text_no_think else [text]

    for src in sources:
        # Try to find a JSON array directly
        for match in re.finditer(r'\[', src):
            start = match.start()
            depth = 0
            for i, ch in enumerate(src[start:], start):
                if ch == '[':
                    depth += 1
                elif ch == ']':
                    depth -= 1
                    if depth == 0:
                        candidate = src[start:i + 1]
                        try:
                            return json.loads(candidate)
                        except json.JSONDecodeError:
                            break

        # Try wrapped in code block: ```json … ```
        cb = re.search(r'```(?:json)?\s*(\[.*?\])\s*```', src, re.DOTALL)
        if cb:
            try:
                return json.loads(cb.group(1))
            except json.JSONDecodeError:
                pass

        # Try a bare JSON object that IS a list item (dict)
        for match in re.finditer(r'\{', src):
            start = match.start()
            depth = 0
            for i, ch in enumerate(src[start:], start):
                if ch == '{':
                    depth += 1
                elif ch == '}':
                    depth -= 1
                    if depth == 0:
                        candidate = src[start:i + 1]
                        try:
                            obj = json.loads(candidate)
                            if isinstance(obj, dict):
                                return [obj]
                        except json.JSONDecodeError:
                            break

    return None


# ── Validation ─────────────────────────────────────────────
VALID_TYPES = {
    "Disease", "Symptom", "Treatment", "Medication",
    "Anatomy", "Provider", "Organization"
}

def _validate_triples(raw_triples: list) -> List[Dict]:
    """Strict quality gate: every triple must have non-empty head/relation/tail."""
    valid = []
    for t in raw_triples:
        if not isinstance(t, dict):
            continue
        head     = str(t.get('head',     '')).strip()
        relation = str(t.get('relation', '')).strip()
        tail     = str(t.get('tail',     '')).strip()
        head_type = str(t.get('head_type', 'Entity')).strip()
        tail_type = str(t.get('tail_type', 'Entity')).strip()

        if len(head) < 2 or len(tail) < 2 or len(relation) < 2:
            continue

        # Normalise types
        if head_type not in VALID_TYPES:
            head_type = 'Entity'
        if tail_type not in VALID_TYPES:
            tail_type = 'Entity'

        valid.append({
            "head":      head,
            "head_type": head_type,
            "relation":  relation,
            "tail":      tail,
            "tail_type": tail_type,
        })
    return valid


# ── Main class ─────────────────────────────────────────────
class KMCKnowledgeGraph:
    def __init__(self):
        logger.info(f"Connecting to Ollama at {OLLAMA_HOST} (GPU 0) …")
        self.client = Client(host=OLLAMA_HOST)

        # Warm-up: ensure model is loaded on GPU 0
        try:
            logger.info(f"Pre-loading model {OLLAMA_MODEL} on GPU 0 …")
            self.client.generate(
                model=OLLAMA_MODEL,
                prompt="Hello",
                options={"num_predict": 1}
            )
            logger.info("Model pre-loaded successfully.")
        except Exception as e:
            logger.warning(f"Model pre-load warning (continuing): {e}")

        try:
            self.driver = GraphDatabase.driver(
                NEO4J_URI, auth=(NEO4J_USER, NEO4J_PASSWORD)
            )
            self.driver.verify_connectivity()
            logger.info("Connected to Neo4j successfully.")
        except Exception as e:
            logger.error(f"Failed to connect to Neo4j: {e}")
            raise

        self.processed_ids = self._load_checkpoints()
        self.lock = threading.RLock()
        
        # Persistent file handles for faster I/O
        self.checkpoint_file = open(CHECKPOINT_PATH, 'a', encoding='utf-8')
        self.output_file = open(OUTPUT_FILE, 'a', encoding='utf-8')
        self.neo4j_ok = True

    # ── Checkpointing ──────────────────────────────────────
    def _load_checkpoints(self) -> set:
        if os.path.exists(CHECKPOINT_PATH):
            with open(CHECKPOINT_PATH, 'r', encoding='utf-8') as f:
                return set(line.strip() for line in f if line.strip())
        return set()

    def _save_checkpoint(self, entry_id: str):
        with self.lock:
            self.checkpoint_file.write(f"{entry_id}\n")
            self.checkpoint_file.flush()
            self.processed_ids.add(entry_id)

    def close(self):
        try:
            self.driver.close()
        except:
            pass
        try:
            self.checkpoint_file.close()
            self.output_file.close()
        except:
            pass

    # ── Text cleaning ──────────────────────────────────────
    def clean_text(self, text: str) -> str:
        if not text:
            return ""
        text = re.sub(r'\s+', ' ', text)
        return text.strip()

    # ── Triple extraction ──────────────────────────────────
    def extract_triples(self, text: str, max_retries: int = 5) -> List[Dict]:
        """
        Extract knowledge-graph triples from Kurdish medical text.

        KEY DESIGN DECISION for Qwen3.6:35b
        ─────────────────────────────────────
        • We do NOT pass format="json" to Ollama.
          Qwen3 uses internal Chain-of-Thought (<think>…</think>) before
          writing the final answer.  Forcing JSON mode suppresses that CoT
          block which drastically hurts extraction quality and often produces
          empty output.
        • Instead we prompt the model to wrap its final JSON in a clear
          delimiter (===JSON_START=== / ===JSON_END===) so our parser can
          find it reliably even inside a long reasoning trace.
        • num_predict is set to 8192 to give the CoT enough room.
        """
        prompt = (
            "You are an elite Kurdish medical linguist and knowledge graph engineer.\n"
            "Your task is to extract high-quality, semantically rich medical triples "
            "from the provided Kurdish text.\n\n"

            "CRITICAL QUALITY REQUIREMENTS:\n"
            "1. LANGUAGE: All medical entities and relationships MUST be in natural, "
            "accurate Kurdish (Sorani/Central Kurdish).\n"
            "2. MEANINGFUL RELATIONS: Use precise medical verbs, never 'has' or 'is'.\n"
            "   Examples:\n"
            "   - 'دەبێتە هۆی'           (causes)\n"
            "   - 'نیشانەیە بۆ'           (is a symptom of)\n"
            "   - 'بەکاردێت بۆ چارەسەری'  (used for treatment of)\n"
            "   - 'کاردەکاتە سەر'         (affects)\n"
            "   - 'بەشێکە لە'            (is part of)\n"
            "3. ENTITY TYPES: head_type and tail_type MUST be ONE of: "
            "[Disease, Symptom, Treatment, Medication, Anatomy, Provider, Organization].\n"
            "4. DATA INTEGRITY: Every triple MUST have a valid head, relation, and tail. "
            "No empty strings.\n"
            "5. COMPLETENESS: Extract ALL meaningful medical facts from the text. "
            "Aim for at least 3-8 triples per passage if the text contains medical info.\n\n"

            f"Kurdish Text:\n{text}\n\n"

            "Think carefully about the medical entities and their relationships.\n"
            "After your reasoning, output YOUR FINAL ANSWER between these exact markers:\n"
            "===JSON_START===\n"
            "[\n"
            "  {\n"
            '    "head": "Entity Name in Kurdish",\n'
            '    "head_type": "English Label",\n'
            '    "relation": "Detailed Kurdish Relationship",\n'
            '    "tail": "Target Entity in Kurdish",\n'
            '    "tail_type": "English Label"\n'
            "  }\n"
            "]\n"
            "===JSON_END===\n\n"
            "If the text contains NO meaningful medical relationships, output:\n"
            "===JSON_START===\n"
            "[]\n"
            "===JSON_END==="
        )

        for attempt in range(max_retries):
            try:
                response = self.client.generate(
                    model=OLLAMA_MODEL,
                    prompt=prompt,
                    # ← NO format="json" — allow full CoT output
                    options=OLLAMA_OPTIONS,
                )

                # Resolve response object
                if isinstance(response, dict):
                    content = response.get('response', '')
                else:
                    content = getattr(response, 'response', '')

                if not content or not content.strip():
                    logger.warning(f"[GPU0] Attempt {attempt+1}: Empty response. Retrying…")
                    time.sleep(1)
                    continue

                logger.debug(f"[GPU0] Raw response length: {len(content)} chars")

                # ── Priority 1: extract between our explicit markers ──
                marker_match = re.search(
                    r'===JSON_START===\s*(.*?)\s*===JSON_END===',
                    content, re.DOTALL
                )
                if marker_match:
                    json_str = marker_match.group(1).strip()
                    try:
                        data = json.loads(json_str)
                        if isinstance(data, list):
                            raw_triples = [i for i in data if isinstance(i, dict)]
                            valid = _validate_triples(raw_triples)
                            if valid:
                                logger.info(f"[GPU0] Attempt {attempt+1}: Extracted {len(valid)} valid triples.")
                                return valid
                            elif not raw_triples:
                                # Model explicitly said no triples
                                logger.info(f"[GPU0] Model found no medical relationships in this text.")
                                return []
                            # else: had triples but all failed validation → retry
                            logger.warning(f"[GPU0] Attempt {attempt+1}: All triples failed validation. Retrying…")
                            continue
                    except json.JSONDecodeError:
                        pass  # Fall through to robust extractor

                # ── Priority 2: robust scan of full response ──
                raw_list = _extract_json_from_text(content)
                if raw_list is not None:
                    if isinstance(raw_list, list):
                        if not raw_list:
                            logger.info(f"[GPU0] Attempt {attempt+1}: Model returned empty list.")
                            return []
                        valid = _validate_triples(raw_list)
                        if valid:
                            logger.info(f"[GPU0] Attempt {attempt+1}: Extracted {len(valid)} valid triples (fallback parser).")
                            return valid
                        logger.warning(f"[GPU0] Attempt {attempt+1}: Fallback triples failed validation. Retrying…")
                        continue

                logger.warning(f"[GPU0] Attempt {attempt+1}: No JSON found in response. Retrying…")
                time.sleep(1)

            except Exception as e:
                logger.warning(f"[GPU0] Attempt {attempt+1} Exception: {e}")
                time.sleep(2)

        logger.error("[GPU0] All retries exhausted. Returning empty list.")
        return []

    # ── Neo4j storage ──────────────────────────────────────
    def store_triples(self, triples: List[Dict]):
        """Store multiple triples in a single Neo4j session for efficiency."""
        if not triples or not self.neo4j_ok:
            return
            
        try:
            with self.driver.session() as session:
                for triple in triples:
                    head     = triple.get('head')
                    h_type   = triple.get('head_type', 'Entity')
                    rel      = triple.get('relation',  'RELATED_TO')
                    tail     = triple.get('tail')
                    t_type   = triple.get('tail_type', 'Entity')

                    if not head or not tail:
                        continue

                    # Neo4j label must be alphanumeric
                    h_label = re.sub(r'[^a-zA-Z0-9]', '', str(h_type)) or "Entity"
                    if h_label[0].isdigit(): h_label = "L_" + h_label

                    t_label = re.sub(r'[^a-zA-Z0-9]', '', str(t_type)) or "Entity"
                    if t_label[0].isdigit(): t_label = "L_" + t_label

                    # Sanitize relation
                    rel_clean = re.sub(r'[^a-zA-Z0-9_]', '_', str(rel)).upper()
                    rel_clean = re.sub(r'_+', '_', rel_clean).strip('_') or "RELATED_TO"
                    if rel_clean[0].isdigit(): rel_clean = "R_" + rel_clean

                    query = (
                        f"MERGE (h:{h_label} {{name: $head}}) "
                        f"MERGE (t:{t_label} {{name: $tail}}) "
                        f"MERGE (h)-[r:{rel_clean}]->(t) "
                    )
                    session.run(query, head=head, tail=tail)
        except Exception as e:
            logger.error(f"[GPU0] Neo4j store error: {e}")
            self.neo4j_ok = False

    # ── Dataset processing ─────────────────────────────────
    def process_entry(self, entry: Dict):
        """Process a single entry from the dataset."""
        entry_id = str(entry.get('id', '')).strip()
        
        with self.lock:
            if not entry_id or entry_id in self.processed_ids:
                return

        raw_text = entry.get('response', '')
        if not raw_text or not raw_text.strip():
            self._save_checkpoint(entry_id)
            return

        cleaned_text = self.clean_text(raw_text)
        triples = self.extract_triples(cleaned_text)

        if triples:
            # ── Persist to JSONL ──
            try:
                record = {"id": entry_id, "triples": triples}
                line = json.dumps(record, ensure_ascii=False) + "\n"
                with self.lock:
                    self.output_file.write(line)
                    self.output_file.flush()
                logger.info(f"[GPU0] ID={entry_id}: Extracted {len(triples)} triples.")
            except Exception as e:
                logger.error(f"[GPU0] Failed to write to {OUTPUT_FILE}: {e}")

            # ── Persist to Neo4j ──
            self.store_triples(triples)

        self._save_checkpoint(entry_id)

    def process_dataset(self):
        logger.info(f"[GPU0] Loading dataset from {DATASET_PATH} …")
        try:
            with open(DATASET_PATH, 'r', encoding='utf-8') as f:
                data = json.load(f)
        except Exception as e:
            logger.error(f"[GPU0] Could not load dataset: {e}")
            return

        total   = len(data)
        with self.lock:
            already = len(self.processed_ids)
        
        remaining_data = [e for e in data if str(e.get('id', '')).strip() not in self.processed_ids]
        logger.info(f"[GPU0] {total} total | {already} done | {len(remaining_data)} remaining.")

        if not remaining_data:
            logger.info("[GPU0] No entries to process.")
            return

        # Use ThreadPoolExecutor for concurrent requests
        # max_workers=3 is a safe balance for GPU 35b models to overlap I/O and compute
        MAX_WORKERS = 3
        with ThreadPoolExecutor(max_workers=MAX_WORKERS) as executor:
            list(tqdm(
                executor.map(self.process_entry, remaining_data), 
                total=len(remaining_data), 
                desc="GPU0"
            ))

        logger.info("[GPU0] Pipeline finished successfully.")


if __name__ == "__main__":
    kg = KMCKnowledgeGraph()
    try:
        kg.process_dataset()
    except KeyboardInterrupt:
        logger.info("[GPU0] Process interrupted by user. Progress saved.")
    except Exception as e:
        logger.error(f"[GPU0] Fatal error: {e}")
    finally:
        kg.close()
\end{lstlisting}

\section{Knowledge Graph Generator (GPU 1)} \label{sec:gengpu1}
\textbf{Purpose:} We use this script to extract medical triples from the second half of the Kurdish dataset using the Ollama service on GPU 1.

\begin{lstlisting}[style=pythoncode, caption={kg\_generator\_gpu\_1.py}]
import json
import os
import re
import time
import logging
import threading
from concurrent.futures import ThreadPoolExecutor
from typing import List, Dict, Optional
from tqdm import tqdm
from ollama import Client
from neo4j import GraphDatabase
from dotenv import load_dotenv

# Load environment variables
load_dotenv()

# ─────────────────────────────────────────────────────────
# GPU 1  →  Ollama on port 11435
# ─────────────────────────────────────────────────────────

# Setup Logging
logging.basicConfig(
    level=logging.INFO,
    format='%(asctime)s - %(levelname)s - %(message)s',
    handlers=[
        logging.FileHandler("kg_pipeline_gpu_1.log"),   # ← correct log file
        logging.StreamHandler()
    ]
)
logger = logging.getLogger(__name__)

# ── Configuration ─────────────────────────────────────────
OLLAMA_HOST   = os.getenv("OLLAMA_HOST_GPU1", "http://localhost:11435")  # ← GPU 1 port
OLLAMA_MODEL  = "qwen3.6:35b"
NEO4J_URI     = os.getenv("NEO4J_URI",      "bolt://localhost:7687")
NEO4J_USER    = os.getenv("NEO4J_USER",     "neo4j")
NEO4J_PASSWORD= os.getenv("NEO4J_PASSWORD", "password")
DATASET_PATH  = "kurdish_medical_corpus_kmc_part2.json"      # ← GPU 1 dataset half
CHECKPOINT_PATH = "processed_ids_gpu_1.txt"                  # ← GPU 1 checkpoint
OUTPUT_FILE   = "extracted_triples_gpu_1.jsonl"              # ← GPU 1 output

# Qwen3 generation options – generous token budgets so CoT has room to breathe
OLLAMA_OPTIONS = {
    "temperature": 0.6,      # Qwen3 docs recommend 0.6 for balanced CoT
    "top_p": 0.95,           # Qwen3 recommended value
    "top_k": 20,             # Qwen3 recommended value
    "min_p": 0,
    "num_predict": 8192,     # Large budget: CoT + JSON output
    "num_ctx": 16384,        # Context window – enough for long Kurdish texts + prompt
    "repeat_penalty": 1.0,
}

# ── Robust JSON extractor ──────────────────────────────────
def _extract_json_from_text(text: str) -> Optional[list]:
    """
    Extract the first JSON array from arbitrary text (including CoT reasoning).
    Qwen3 thinks in <think>…</think> blocks before giving the answer;
    we must scan the *entire* response for any valid JSON array.
    """
    if not text:
        return None

    # Strip think blocks if present (Qwen3 CoT)
    text_no_think = re.sub(r'<think>.*?</think>', '', text, flags=re.DOTALL).strip()

    # Prefer text after think block; fall back to full text
    sources = [text_no_think, text] if text_no_think else [text]

    for src in sources:
        # Try to find a JSON array directly
        for match in re.finditer(r'\[', src):
            start = match.start()
            depth = 0
            for i, ch in enumerate(src[start:], start):
                if ch == '[':
                    depth += 1
                elif ch == ']':
                    depth -= 1
                    if depth == 0:
                        candidate = src[start:i + 1]
                        try:
                            return json.loads(candidate)
                        except json.JSONDecodeError:
                            break

        # Try wrapped in code block: ```json … ```
        cb = re.search(r'```(?:json)?\s*(\[.*?\])\s*```', src, re.DOTALL)
        if cb:
            try:
                return json.loads(cb.group(1))
            except json.JSONDecodeError:
                pass

        # Try a bare JSON object that IS a list item (dict)
        for match in re.finditer(r'\{', src):
            start = match.start()
            depth = 0
            for i, ch in enumerate(src[start:], start):
                if ch == '{':
                    depth += 1
                elif ch == '}':
                    depth -= 1
                    if depth == 0:
                        candidate = src[start:i + 1]
                        try:
                            obj = json.loads(candidate)
                            if isinstance(obj, dict):
                                return [obj]
                        except json.JSONDecodeError:
                            break

    return None


# ── Validation ─────────────────────────────────────────────
VALID_TYPES = {
    "Disease", "Symptom", "Treatment", "Medication",
    "Anatomy", "Provider", "Organization"
}

def _validate_triples(raw_triples: list) -> List[Dict]:
    """Strict quality gate: every triple must have non-empty head/relation/tail."""
    valid = []
    for t in raw_triples:
        if not isinstance(t, dict):
            continue
        head      = str(t.get('head',      '')).strip()
        relation  = str(t.get('relation',  '')).strip()
        tail      = str(t.get('tail',      '')).strip()
        head_type = str(t.get('head_type', 'Entity')).strip()
        tail_type = str(t.get('tail_type', 'Entity')).strip()

        if len(head) < 2 or len(tail) < 2 or len(relation) < 2:
            continue

        # Normalise types
        if head_type not in VALID_TYPES:
            head_type = 'Entity'
        if tail_type not in VALID_TYPES:
            tail_type = 'Entity'

        valid.append({
            "head":      head,
            "head_type": head_type,
            "relation":  relation,
            "tail":      tail,
            "tail_type": tail_type,
        })
    return valid


# ── Main class ─────────────────────────────────────────────
class KMCKnowledgeGraph:
    def __init__(self):
        logger.info(f"Connecting to Ollama at {OLLAMA_HOST} (GPU 1) …")
        self.client = Client(host=OLLAMA_HOST)

        # Warm-up: ensure model is loaded on GPU 1
        try:
            logger.info(f"Pre-loading model {OLLAMA_MODEL} on GPU 1 …")
            self.client.generate(
                model=OLLAMA_MODEL,
                prompt="Hello",
                options={"num_predict": 1}
            )
            logger.info("Model pre-loaded successfully.")
        except Exception as e:
            logger.warning(f"Model pre-load warning (continuing): {e}")

        try:
            self.driver = GraphDatabase.driver(
                NEO4J_URI, auth=(NEO4J_USER, NEO4J_PASSWORD)
            )
            self.driver.verify_connectivity()
            logger.info("Connected to Neo4j successfully.")
        except Exception as e:
            logger.error(f"Failed to connect to Neo4j: {e}")
            raise

        self.processed_ids = self._load_checkpoints()
        self.lock = threading.RLock()
        
        # Persistent file handles for faster I/O
        self.checkpoint_file = open(CHECKPOINT_PATH, 'a', encoding='utf-8')
        self.output_file = open(OUTPUT_FILE, 'a', encoding='utf-8')
        self.neo4j_ok = True

    # ── Checkpointing ──────────────────────────────────────
    def _load_checkpoints(self) -> set:
        if os.path.exists(CHECKPOINT_PATH):
            with open(CHECKPOINT_PATH, 'r', encoding='utf-8') as f:
                return set(line.strip() for line in f if line.strip())
        return set()

    def _save_checkpoint(self, entry_id: str):
        with self.lock:
            self.checkpoint_file.write(f"{entry_id}\n")
            self.checkpoint_file.flush()
            self.processed_ids.add(entry_id)

    def close(self):
        try:
            self.driver.close()
        except:
            pass
        try:
            self.checkpoint_file.close()
            self.output_file.close()
        except:
            pass

    # ── Text cleaning ──────────────────────────────────────
    def clean_text(self, text: str) -> str:
        if not text:
            return ""
        text = re.sub(r'\s+', ' ', text)
        return text.strip()

    # ── Triple extraction ──────────────────────────────────
    def extract_triples(self, text: str, max_retries: int = 5) -> List[Dict]:
        """
        Extract knowledge-graph triples from Kurdish medical text.

        KEY DESIGN DECISION for Qwen3.6:35b
        ─────────────────────────────────────
        • We do NOT pass format="json" to Ollama.
          Qwen3 uses internal Chain-of-Thought (<think>…</think>) before
          writing the final answer.  Forcing JSON mode suppresses that CoT
          block which drastically hurts extraction quality and often produces
          empty output.
        • Instead we prompt the model to wrap its final JSON in a clear
          delimiter (===JSON_START=== / ===JSON_END===) so our parser can
          find it reliably even inside a long reasoning trace.
        • num_predict is set to 8192 to give the CoT enough room.
        """
        prompt = (
            "You are an elite Kurdish medical linguist and knowledge graph engineer.\n"
            "Your task is to extract high-quality, semantically rich medical triples "
            "from the provided Kurdish text.\n\n"

            "CRITICAL QUALITY REQUIREMENTS:\n"
            "1. LANGUAGE: All medical entities and relationships MUST be in natural, "
            "accurate Kurdish (Sorani/Central Kurdish).\n"
            "2. MEANINGFUL RELATIONS: Use precise medical verbs, never 'has' or 'is'.\n"
            "   Examples:\n"
            "   - 'دەبێتە هۆی'           (causes)\n"
            "   - 'نیشانەیە بۆ'           (is a symptom of)\n"
            "   - 'بەکاردێت بۆ چارەسەری'  (used for treatment of)\n"
            "   - 'کاردەکاتە سەر'         (affects)\n"
            "   - 'بەشێکە لە'            (is part of)\n"
            "3. ENTITY TYPES: head_type and tail_type MUST be ONE of: "
            "[Disease, Symptom, Treatment, Medication, Anatomy, Provider, Organization].\n"
            "4. DATA INTEGRITY: Every triple MUST have a valid head, relation, and tail. "
            "No empty strings.\n"
            "5. COMPLETENESS: Extract ALL meaningful medical facts from the text. "
            "Aim for at least 3-8 triples per passage if the text contains medical info.\n\n"

            f"Kurdish Text:\n{text}\n\n"

            "Think carefully about the medical entities and their relationships.\n"
            "After your reasoning, output YOUR FINAL ANSWER between these exact markers:\n"
            "===JSON_START===\n"
            "[\n"
            "  {\n"
            '    "head": "Entity Name in Kurdish",\n'
            '    "head_type": "English Label",\n'
            '    "relation": "Detailed Kurdish Relationship",\n'
            '    "tail": "Target Entity in Kurdish",\n'
            '    "tail_type": "English Label"\n'
            "  }\n"
            "]\n"
            "===JSON_END===\n\n"
            "If the text contains NO meaningful medical relationships, output:\n"
            "===JSON_START===\n"
            "[]\n"
            "===JSON_END==="
        )

        for attempt in range(max_retries):
            try:
                response = self.client.generate(
                    model=OLLAMA_MODEL,
                    prompt=prompt,
                    # ← NO format="json" — allow full CoT output
                    options=OLLAMA_OPTIONS,
                )

                # Resolve response object
                if isinstance(response, dict):
                    content = response.get('response', '')
                else:
                    content = getattr(response, 'response', '')

                if not content or not content.strip():
                    logger.warning(f"[GPU1] Attempt {attempt+1}: Empty response. Retrying…")
                    time.sleep(1)
                    continue

                logger.debug(f"[GPU1] Raw response length: {len(content)} chars")

                # ── Priority 1: extract between our explicit markers ──
                marker_match = re.search(
                    r'===JSON_START===\s*(.*?)\s*===JSON_END===',
                    content, re.DOTALL
                )
                if marker_match:
                    json_str = marker_match.group(1).strip()
                    try:
                        data = json.loads(json_str)
                        if isinstance(data, list):
                            raw_triples = [i for i in data if isinstance(i, dict)]
                            valid = _validate_triples(raw_triples)
                            if valid:
                                logger.info(f"[GPU1] Attempt {attempt+1}: Extracted {len(valid)} valid triples.")
                                return valid
                            elif not raw_triples:
                                # Model explicitly said no triples
                                logger.info(f"[GPU1] Model found no medical relationships in this text.")
                                return []
                            # else: had triples but all failed validation → retry
                            logger.warning(f"[GPU1] Attempt {attempt+1}: All triples failed validation. Retrying…")
                            continue
                    except json.JSONDecodeError:
                        pass  # Fall through to robust extractor

                # ── Priority 2: robust scan of full response ──
                raw_list = _extract_json_from_text(content)
                if raw_list is not None:
                    if isinstance(raw_list, list):
                        if not raw_list:
                            logger.info(f"[GPU1] Attempt {attempt+1}: Model returned empty list.")
                            return []
                        valid = _validate_triples(raw_list)
                        if valid:
                            logger.info(f"[GPU1] Attempt {attempt+1}: Extracted {len(valid)} valid triples (fallback parser).")
                            return valid
                        logger.warning(f"[GPU1] Attempt {attempt+1}: Fallback triples failed validation. Retrying…")
                        continue

                logger.warning(f"[GPU1] Attempt {attempt+1}: No JSON found in response. Retrying…")
                time.sleep(1)

            except Exception as e:
                logger.warning(f"[GPU1] Attempt {attempt+1} Exception: {e}")
                time.sleep(2)

        logger.error("[GPU1] All retries exhausted. Returning empty list.")
        return []

    # ── Neo4j storage ──────────────────────────────────────
    def store_triples(self, triples: List[Dict]):
        """Store multiple triples in a single Neo4j session for efficiency."""
        if not triples or not self.neo4j_ok:
            return
            
        try:
            with self.driver.session() as session:
                for triple in triples:
                    head     = triple.get('head')
                    h_type   = triple.get('head_type', 'Entity')
                    rel      = triple.get('relation',  'RELATED_TO')
                    tail     = triple.get('tail')
                    t_type   = triple.get('tail_type', 'Entity')

                    if not head or not tail:
                        continue

                    # Neo4j label must be alphanumeric
                    h_label = re.sub(r'[^a-zA-Z0-9]', '', str(h_type)) or "Entity"
                    if h_label[0].isdigit(): h_label = "L_" + h_label

                    t_label = re.sub(r'[^a-zA-Z0-9]', '', str(t_type)) or "Entity"
                    if t_label[0].isdigit(): t_label = "L_" + t_label

                    # Sanitize relation
                    rel_clean = re.sub(r'[^a-zA-Z0-9_]', '_', str(rel)).upper()
                    rel_clean = re.sub(r'_+', '_', rel_clean).strip('_') or "RELATED_TO"
                    if rel_clean[0].isdigit(): rel_clean = "R_" + rel_clean

                    query = (
                        f"MERGE (h:{h_label} {{name: $head}}) "
                        f"MERGE (t:{t_label} {{name: $tail}}) "
                        f"MERGE (h)-[r:{rel_clean}]->(t) "
                    )
                    session.run(query, head=head, tail=tail)
        except Exception as e:
            logger.error(f"[GPU1] Neo4j store error: {e}")
            self.neo4j_ok = False

    # ── Dataset processing ─────────────────────────────────
    def process_entry(self, entry: Dict):
        """Process a single entry from the dataset."""
        entry_id = str(entry.get('id', '')).strip()
        
        with self.lock:
            if not entry_id or entry_id in self.processed_ids:
                return

        raw_text = entry.get('response', '')
        if not raw_text or not raw_text.strip():
            self._save_checkpoint(entry_id)
            return

        cleaned_text = self.clean_text(raw_text)
        triples = self.extract_triples(cleaned_text)

        if triples:
            # ── Persist to JSONL ──
            try:
                record = {"id": entry_id, "triples": triples}
                line = json.dumps(record, ensure_ascii=False) + "\n"
                with self.lock:
                    self.output_file.write(line)
                    self.output_file.flush()
                logger.info(f"[GPU1] ID={entry_id}: Extracted {len(triples)} triples.")
            except Exception as e:
                logger.error(f"[GPU1] Failed to write to {OUTPUT_FILE}: {e}")

            # ── Persist to Neo4j ──
            self.store_triples(triples)

        self._save_checkpoint(entry_id)

    def process_dataset(self):
        logger.info(f"[GPU1] Loading dataset from {DATASET_PATH} …")
        try:
            with open(DATASET_PATH, 'r', encoding='utf-8') as f:
                data = json.load(f)
        except Exception as e:
            logger.error(f"[GPU1] Could not load dataset: {e}")
            return

        total   = len(data)
        with self.lock:
            already = len(self.processed_ids)
        
        remaining_data = [e for e in data if str(e.get('id', '')).strip() not in self.processed_ids]
        logger.info(f"[GPU1] {total} total | {already} done | {len(remaining_data)} remaining.")

        if not remaining_data:
            logger.info("[GPU1] No entries to process.")
            return

        # Use ThreadPoolExecutor for concurrent requests
        # max_workers=3 is a safe balance for GPU 35b models to overlap I/O and compute
        MAX_WORKERS = 3
        with ThreadPoolExecutor(max_workers=MAX_WORKERS) as executor:
            list(tqdm(
                executor.map(self.process_entry, remaining_data), 
                total=len(remaining_data), 
                desc="GPU1"
            ))

        logger.info("[GPU1] Pipeline finished successfully.")


if __name__ == "__main__":
    kg = KMCKnowledgeGraph()
    try:
        kg.process_dataset()
    except KeyboardInterrupt:
        logger.info("[GPU1] Process interrupted by user. Progress saved.")
    except Exception as e:
        logger.error(f"[GPU1] Fatal error: {e}")
    finally:
        kg.close()
\end{lstlisting}

\section{Run KG Script (GPU 0)} \label{sec:runkg0}
\textbf{Purpose:} A shell script to execute the GPU 0 knowledge graph generator in the background.

\begin{lstlisting}[style=bashcode, caption={run\_kg\_gpu\_0.sh}]
#!/bin/bash

# Check if virtual environment exists and activate it
if [ -d ".venv" ]; then
    echo "Activating virtual environment..."
    source .venv/bin/activate
elif [ -d "venv" ]; then
    echo "Activating virtual environment (venv)..."
    source venv/bin/activate
fi

# Run the kg_generator_gpu_0.py in the background
# We use -u for unbuffered output to ensure logs appear in real-time
export PYTHONUNBUFFERED=1
nohup python3 -u kg_generator_gpu_0.py > kg_console_gpu_0.log 2>&1 &

PID=$!
echo "--------------------------------------------------"
echo "KG Generator (GPU 0) started in background."
echo "PID: $PID"
echo "Console output: kg_console_gpu_0.log"
echo "To stop the process, run: kill $PID"
echo "--------------------------------------------------"
\end{lstlisting}

\section{Run KG Script (GPU 1)} \label{sec:runkg1}
\textbf{Purpose:} A shell script to execute the GPU 1 knowledge graph generator in the background.

\begin{lstlisting}[style=bashcode, caption={run\_kg\_gpu\_1.sh}]
#!/bin/bash

# Check if virtual environment exists and activate it
if [ -d ".venv" ]; then
    echo "Activating virtual environment..."
    source .venv/bin/activate
elif [ -d "venv" ]; then
    echo "Activating virtual environment (venv)..."
    source venv/bin/activate
fi

# Run the kg_generator_gpu_1.py in the background
# We use -u for unbuffered output to ensure logs appear in real-time
export PYTHONUNBUFFERED=1
nohup python3 -u kg_generator_gpu_1.py > kg_console_gpu_1.log 2>&1 &

PID=$!
echo "--------------------------------------------------"
echo "KG Generator (GPU 1) started in background."
echo "PID: $PID"
echo "Console output: kg_console_gpu_1.log"
echo "To stop the process, run: kill $PID"
echo "--------------------------------------------------"
\end{lstlisting}

\section{Stop KG Script} \label{sec:stopkg}
\textbf{Purpose:} Gracefully stops the background Python generator processes and terminates any stray Ollama server instances.

\begin{lstlisting}[style=bashcode, caption={stop\_kg.sh}]
#!/bin/bash

# ==============================================================================
# KMC Knowledge Graph Pipeline Shutdown Script
# ==============================================================================

echo "--------------------------------------------------"
echo "🛑 Initiating shutdown of KG Pipeline and Models..."
echo "--------------------------------------------------"

# 1. Stop Python Generators
stop_python_process() {
    local script_name=$1
    local gpu_id=$2
    local pid=$(pgrep -f "python3.*$script_name")
    
    if [ -n "$pid" ]; then
        echo "[*] GPU $gpu_id: Stopping $script_name (PID: $pid)..."
        kill $pid
        sleep 2
        # Force kill if still running
        if ps -p $pid > /dev/null; then
            echo "[!] GPU $gpu_id: Process persistent. Force killing..."
            kill -9 $pid
        fi
    else
        echo "[-] GPU $gpu_id: No generator process found for $script_name."
    fi
}

stop_python_process "kg_generator_gpu_0.py" 0
stop_python_process "kg_generator_gpu_1.py" 1

# 2. Stop Ollama Instances
echo "[*] Cleaning up Ollama server and model instances..."
if pgrep -f "ollama" > /dev/null; then
    echo "[!] Found active Ollama processes. Terminating..."
    pkill -f "ollama"
    sleep 3
    # Check again and force kill if necessary
    if pgrep -f "ollama" > /dev/null; then
        echo "[!] Some Ollama processes are stuck. Force killing all..."
        pkill -9 -f "ollama"
    fi
    echo "[+] Ollama processes terminated."
else
    echo "[-] No Ollama processes found."
fi

echo "--------------------------------------------------"
echo "✅ Shutdown complete."
\end{lstlisting}

\section{Setup Script} \label{sec:setup}
\textbf{Purpose:} We use this shell script to prepare the environment, install required Python dependencies from \texttt{requirements.txt}, initialize a Neo4j Docker container, and verify the Ollama installation and model availability.

\begin{lstlisting}[style=bashcode, caption={setup.sh}]
#!/bin/bash

# KMC Knowledge Graph Setup Script

echo "🚀 Starting setup for KMC Knowledge Graph Pipeline..."

# 1. Check for Python
if ! command -v python3 &> /dev/null; then
    echo "❌ Python3 is not installed. Please install it first."
    exit 1
fi

# 2. Install Requirements
echo "📥 Installing dependencies from requirements.txt..."
pip install --upgrade pip
pip install -r requirements.txt

# 3. Docker and Neo4j Setup
echo "🐳 Checking Docker status..."
if ! command -v docker &> /dev/null; then
    echo "📥 Docker not found. Installing Docker..."
    curl -fsSL https://get.docker.com -o get-docker.sh
    sudo sh get-docker.sh
    sudo usermod -aG docker $USER
    rm get-docker.sh
    DOCKER_CMD="sudo docker"
    echo "✅ Docker installed. Using 'sudo docker' for this session."
else
    DOCKER_CMD="docker"
fi

echo "🚢 Setting up Neo4j container..."
if ! $DOCKER_CMD ps -a --format '{{.Names}}' | grep -q "^neo4j-kmc$"; then
    $DOCKER_CMD run -d \
        --name neo4j-kmc \
        -p 7474:7474 -p 7687:7687 \
        -e NEO4J_AUTH=neo4j/password \
        -e NEO4J_PLUGINS='["apoc"]' \
        --restart always \
        neo4j:latest
    echo "✅ Neo4j started. Access it at http://localhost:7474"
else
    echo "✅ Neo4j container 'neo4j-kmc' already exists."
    if [ "$($DOCKER_CMD inspect -f '{{.State.Running}}' neo4j-kmc)" != "true" ]; then
        echo "🔄 Starting Neo4j container..."
        $DOCKER_CMD start neo4j-kmc
    fi
fi

# 4. Check for Ollama
if ! command -v ollama &> /dev/null; then
    echo "⚠️ Ollama not found in PATH. Please install it from https://ollama.com"
else
    if ollama list | grep -q "qwen3.6:35b"; then
        echo "✅ Model qwen3.6:35b already exists in Ollama."
    else
        echo "🧠 Pulling Qwen model (qwen3.6:35b)..."
        ollama pull qwen3.6:35b
    fi
fi

# 5. Create .env template if it doesn't exist
if [ ! -f .env ]; then
    echo "📝 Creating .env template..."
    echo "NEO4J_URI=bolt://localhost:7687" > .env
    echo "NEO4J_USER=neo4j" >> .env
    echo "NEO4J_PASSWORD=password" >> .env
fi

echo "✅ Setup complete!"
echo "💡 You can now run the pipeline: 'python kg_generator.py'"
\end{lstlisting}

\section{Setup Runner} \label{sec:runsetup}
\textbf{Purpose:} A Python utility script to launch the \texttt{setup.sh} installation routine in a detached background process.

\begin{lstlisting}[style=pythoncode, caption={run\_setup.py}]
import subprocess
import os
import sys

def run_setup_background():
    """
    Runs setup.sh in the background and redirects all output to setup.log.
    Designed for execution on an Ubuntu server.
    """
    script_name = "setup.sh"
    log_name = "setup.log"
    
    # Check if setup.sh exists
    if not os.path.exists(script_name):
        print(f"❌ Error: {script_name} not found in the current directory.")
        return

    # Ensure the script is executable
    os.chmod(script_name, 0o755)

    print(f"🚀 Launching {script_name} in the background...")
    print(f"📋 Output will be logged to {log_name}")

    try:
        # Open the log file for writing
        with open(log_name, "w") as log_file:
            # subprocess.Popen starts the process without waiting for it to finish
            # We use 'bash' explicitly to ensure it runs correctly on Ubuntu
            process = subprocess.Popen(
                ["bash", script_name],
                stdout=log_file,
                stderr=subprocess.STDOUT, # Redirect stderr to the same log file
                start_new_session=True    # Detach from the current terminal session
            )
        
        print(f"✅ Setup background process started successfully!")
        print(f"🆔 PID: {process.pid}")
        print(f"💡 You can monitor progress with: tail -f {log_name}")

    except Exception as e:
        print(f"❌ Failed to start the setup process: {e}")

if __name__ == "__main__":
    run_setup_background()
\end{lstlisting}

\section{Model Loader (GPU 0)} \label{sec:loader0}
\textbf{Purpose:} A sophisticated bash script to trigger the loading of the Qwen model into GPU 0 memory in the background and verify its active status via polling.

\begin{lstlisting}[style=bashcode, caption={load\_model\_gpu\_0.sh}]
#!/bin/bash

# ==============================================================================
# Ollama Model GPU Loader & Verification Script (GPU 0)
# ==============================================================================
# Description: This script triggers the loading of an Ollama model into GPU 0
#              memory in the background and verifies the status.
# ==============================================================================

# --- Configuration ---
MODEL_NAME="qwen3.6:35b"
LOG_FILE="ollama_gpu_0_load.txt"
GPU_ID=0  # The index of the GPU to use
OLLAMA_PORT=11434

# ANSI Color Codes for Premium UI
GREEN='\033[0;32m'
BLUE='\033[0;34m'
CYAN='\033[0;36m'
YELLOW='\033[1;33m'
RED='\033[0;31m'
BOLD='\033[1m'
NC='\033[0m' # No Color

# Header
echo -e "${CYAN}${BOLD}======================================================${NC}"
echo -e "${CYAN}${BOLD}       OLLAMA GPU 0 LOAD & TEST UTILITY               ${NC}"
echo -e "${CYAN}${BOLD}======================================================${NC}"

# Initialize log file
echo "--- Ollama GPU 0 Load Session: $(date) ---" > "$LOG_FILE"

# --- Function: Ensure Server is Running ---
ensure_server_running() {
    echo -e "${YELLOW}[*] Action: Checking Ollama server status on port $OLLAMA_PORT...${NC}"
    
    if curl -s "http://localhost:$OLLAMA_PORT/api/tags" > /dev/null 2>&1; then
        echo -e "${GREEN}[+] Ollama server is already active on port $OLLAMA_PORT.${NC}"
        return 0
    fi

    echo -e "${YELLOW}[!] Ollama server not detected on port $OLLAMA_PORT. Starting isolated instance...${NC}"
    
    # CRITICAL: Set CUDA_VISIBLE_DEVICES for the SERVER process
    export CUDA_VISIBLE_DEVICES=$GPU_ID
    export OLLAMA_HOST="127.0.0.1:$OLLAMA_PORT"
    
    # Start server in background
    nohup ollama serve >> "$LOG_FILE" 2>&1 &

    # Wait and verify status
    echo -n "Waiting for API availability"
    for i in {1..20}; do
        echo -n "."
        if curl -s "http://localhost:$OLLAMA_PORT/api/tags" > /dev/null 2>&1; then
            echo -e "\n${GREEN}[SUCCESS] Ollama server (GPU $GPU_ID) is up!${NC}" | tee -a "$LOG_FILE"
            return 0
        fi
        sleep 2
    done

    echo -e "\n${RED}[ERROR] Server failed to start within 40 seconds.${NC}" | tee -a "$LOG_FILE"
    exit 1
}

# --- Function: Load Model ---
load_to_gpu_bg() {
    echo -e "${YELLOW}[*] Action: Triggering model load for '$MODEL_NAME' on port $OLLAMA_PORT...${NC}"
    
    # Trigger a run with empty input to force the model into VRAM
    export OLLAMA_HOST="127.0.0.1:$OLLAMA_PORT"
    nohup ollama run $MODEL_NAME "" >> "$LOG_FILE" 2>&1 &
    
    BG_PID=$!
    echo -e "${GREEN}[+] Load command dispatched (PID: $BG_PID).${NC}" | tee -a "$LOG_FILE"
}

# --- Function: Verification Test ---
run_verification_test() {
    echo -e "\n${YELLOW}[*] Action: Verifying GPU VRAM status...${NC}"
    echo -n "Waiting for model initialization"
    
    MAX_ATTEMPTS=12
    SUCCESS=false
    
    for ((i=1; i<=MAX_ATTEMPTS; i++)); do
        echo -n "."
        if OLLAMA_HOST="localhost:$OLLAMA_PORT" ollama ps 2>/dev/null | grep -q "$MODEL_NAME"; then
            echo -e "\n${GREEN}${BOLD}[SUCCESS] Model '$MODEL_NAME' is active in GPU memory!${NC}" | tee -a "$LOG_FILE"
            SUCCESS=true
            break
        fi
        sleep 5
    done

    if [ "$SUCCESS" = true ]; then
        echo -e "${BLUE}------------------------------------------------------${NC}"
        echo -e "${BOLD}Current GPU Load Status (Port $OLLAMA_PORT):${NC}"
        OLLAMA_HOST="localhost:$OLLAMA_PORT" ollama ps | grep -E "NAME|${MODEL_NAME}"
        echo -e "${BLUE}------------------------------------------------------${NC}"
        return 0
    else
        echo -e "\n${RED}${BOLD}[FAILURE] Model verification timed out.${NC}" | tee -a "$LOG_FILE"
        return 1
    fi
}

# --- Main Execution ---
ensure_server_running
load_to_gpu_bg
run_verification_test

echo -e "\n${CYAN}Done.${NC}"
\end{lstlisting}

\section{Model Loader (GPU 1)} \label{sec:loader1}
\textbf{Purpose:} A sophisticated bash script to trigger the loading of the Qwen model into GPU 1 memory in the background and verify its active status via polling.

\begin{lstlisting}[style=bashcode, caption={load\_model\_gpu\_1.sh}]
#!/bin/bash

# ==============================================================================
# Ollama Model GPU Loader & Verification Script (GPU 1)
# ==============================================================================
# Description: This script triggers the loading of an Ollama model into GPU 1
#              memory in the background and verifies the status.
# ==============================================================================

# --- Configuration ---
MODEL_NAME="qwen3.6:35b"
LOG_FILE="ollama_gpu_1_load.txt"
GPU_ID=1  # The index of the GPU to use
OLLAMA_PORT=11435

# ANSI Color Codes for Premium UI
GREEN='\033[0;32m'
BLUE='\033[0;34m'
CYAN='\033[0;36m'
YELLOW='\033[1;33m'
RED='\033[0;31m'
BOLD='\033[1m'
NC='\033[0m' # No Color

# Header
echo -e "${CYAN}${BOLD}======================================================${NC}"
echo -e "${CYAN}${BOLD}       OLLAMA GPU 1 LOAD & TEST UTILITY               ${NC}"
echo -e "${CYAN}${BOLD}======================================================${NC}"

# Initialize log file
echo "--- Ollama GPU 1 Load Session: $(date) ---" > "$LOG_FILE"

# --- Function: Ensure Server is Running ---
ensure_server_running() {
    echo -e "${YELLOW}[*] Action: Checking Ollama server status on port $OLLAMA_PORT...${NC}"
    
    if curl -s "http://localhost:$OLLAMA_PORT/api/tags" > /dev/null 2>&1; then
        echo -e "${GREEN}[+] Ollama server is already active on port $OLLAMA_PORT.${NC}"
        return 0
    fi

    echo -e "${YELLOW}[!] Ollama server not detected on port $OLLAMA_PORT. Starting isolated instance...${NC}"
    
    # CRITICAL: Set CUDA_VISIBLE_DEVICES for the SERVER process
    export CUDA_VISIBLE_DEVICES=$GPU_ID
    export OLLAMA_HOST="127.0.0.1:$OLLAMA_PORT"
    
    # Start server in background
    nohup ollama serve >> "$LOG_FILE" 2>&1 &

    # Wait and verify status
    echo -n "Waiting for API availability"
    for i in {1..20}; do
        echo -n "."
        if curl -s "http://localhost:$OLLAMA_PORT/api/tags" > /dev/null 2>&1; then
            echo -e "\n${GREEN}[SUCCESS] Ollama server (GPU $GPU_ID) is up!${NC}" | tee -a "$LOG_FILE"
            return 0
        fi
        sleep 2
    done

    echo -e "\n${RED}[ERROR] Server failed to start within 40 seconds.${NC}" | tee -a "$LOG_FILE"
    exit 1
}

# --- Function: Load Model ---
load_to_gpu_bg() {
    echo -e "${YELLOW}[*] Action: Triggering model load for '$MODEL_NAME' on port $OLLAMA_PORT...${NC}"
    
    # Trigger a run with empty input to force the model into VRAM
    export OLLAMA_HOST="127.0.0.1:$OLLAMA_PORT"
    nohup ollama run $MODEL_NAME "" >> "$LOG_FILE" 2>&1 &
    
    BG_PID=$!
    echo -e "${GREEN}[+] Load command dispatched (PID: $BG_PID).${NC}" | tee -a "$LOG_FILE"
}

# --- Function: Verification Test ---
run_verification_test() {
    echo -e "\n${YELLOW}[*] Action: Verifying GPU VRAM status...${NC}"
    echo -n "Waiting for model initialization"
    
    MAX_ATTEMPTS=12
    SUCCESS=false
    
    for ((i=1; i<=MAX_ATTEMPTS; i++)); do
        echo -n "."
        if OLLAMA_HOST="localhost:$OLLAMA_PORT" ollama ps 2>/dev/null | grep -q "$MODEL_NAME"; then
            echo -e "\n${GREEN}${BOLD}[SUCCESS] Model '$MODEL_NAME' is active in GPU memory!${NC}" | tee -a "$LOG_FILE"
            SUCCESS=true
            break
        fi
        sleep 5
    done

    if [ "$SUCCESS" = true ]; then
        echo -e "${BLUE}------------------------------------------------------${NC}"
        echo -e "${BOLD}Current GPU Load Status (Port $OLLAMA_PORT):${NC}"
        OLLAMA_HOST="localhost:$OLLAMA_PORT" ollama ps | grep -E "NAME|${MODEL_NAME}"
        echo -e "${BLUE}------------------------------------------------------${NC}"
        return 0
    else
        echo -e "\n${RED}${BOLD}[FAILURE] Model verification timed out.${NC}" | tee -a "$LOG_FILE"
        return 1
    fi
}

# --- Main Execution ---
ensure_server_running
load_to_gpu_bg
run_verification_test

echo -e "\n${CYAN}Done.${NC}"
\end{lstlisting}

\section{Requirements} \label{sec:requirements}
\textbf{Purpose:} Defines the required Python packages for running the dual-GPU knowledge graph pipeline.

\begin{lstlisting}[style=bashcode, caption={requirements.txt}]
neo4j
ollama
python-dotenv
tqdm
\end{lstlisting}

\end{document}
